\medskip

\textbf{Les questions 1, 2, 3 et 4 sont indépendantes les unes des autres.}

\bigskip

\textbf{Question 1}

\medskip

\emph{Pour cette question, indiquer, en justifiant, la lettre correspondant à la réponse exacte.}

\medskip

On considère la fonction$f$ définie sur $\R$ par $f(x)=(2 + 5x)\e^{3x}$.
On admet que $f$ est dérivable sur $\R$ et on note $f'$ sa dérivée.
Pour tout $x$ appartenant à $\R$, on a :

\renewcommand\arraystretch{1.6}
\begin{tabularx}{\linewidth}{|*{4}{>{\centering \arraybackslash}X|}}\hline
A&B&C&D\\\hline
$f'(x) = 5\e^{3x}$&$f'(x) = 15\e^{3x}$&
$f'(x) = (11 + 15x)\e^{3x}$&$f'(x) = (7 + 5x)\e^{2x}$\\\hline
\end{tabularx}
\renewcommand\arraystretch{1}

\bigskip

\textbf{Question 2}

\medskip

On considère l'équation différentielle $(E)$ :

\[y' = - 3y + 5,\]

où $y$ est une fonction de la variable $x$, définie et dérivable sur $\R$.

Déterminer les fonctions définies sur $\R$, solutions de l'équation différentielle $(E)$.

\bigskip

\textbf{Question 3}

\medskip

Déterminer la forme exponentielle du nombre complexe :
\[z = -6\sqrt{3} + 6\i.\]

\bigskip

\textbf{Question 4}

\medskip

\begin{minipage}{0.6\linewidth}
Les fonctions $f$ et $g$ définies sur $\R$ respectivement par :

\[f(x) =\dfrac12 x^2 +1 \text{ et } g(x) = x + 5,\]

sont représentées sur le graphique ci-contre par la courbe $\mathcal{C}$, courbe représentative de la fonction $f$ et la droite $\Delta$, courbe représentative de la fonction $g$.

Lire graphiquement les positions relatives des courbes représentatives $\mathcal{C}$ de la fonction $f$ et $\Delta$ de la fonction $g$ puis montrer que l'aire de la partie colorée comprise entre la courbe $\mathcal{C}$ et la droite $\Delta$ vaut 18 unités d'aire.
\end{minipage}\hfill
\begin{minipage}{0.39\linewidth}
\psset{unit=0.5cm,arrowsize=2pt 3}
\begin{pspicture*}(-6,-1)(5.5,10)
\pscustom[fillstyle=solid,fillcolor=lightgray]{\psplot[plotpoints=2000,linewidth=1.25pt,linecolor=red]{-2}{4}{x dup mul 0.5 mul 1 add}\psline(4,9)(-2,3)
}
\psgrid[gridlabels=0pt,subgriddiv=1,gridwidth=0.25pt]
\psaxes[linewidth=1.25pt,labelFontSize=\scriptstyle]{->}(0,0)(-6,-1)(5.5,10)
\psplot[plotpoints=2000,linewidth=1.25pt,linecolor=red]{-6}{5.5}{x dup mul 0.5 mul 1 add}
\psplot[plotpoints=500,linewidth=1.25pt,linecolor=blue]{-6}{5.5}{x 5 add}
\uput[r](-4,9){\red $\mathcal{C}$}\uput[ul](2,7){\blue $\Delta$}
\end{pspicture*}
\end{minipage}

\bigskip
